%% ============================================================================
%% Elsevier CAS single-column manuscript
%% Three diagnostics for trustworthy claims in neural-coordinate WE sampling
%%
%% Build: pdflatex x2 (no external .bib; thebibliography is inline)
%% Requires: cas-sc.cls (Elsevier CAS bundle)
%% ============================================================================

\documentclass[a4paper,fleqn]{cas-sc}

\usepackage[numbers,sort&compress]{natbib}
\usepackage{amsmath}
\usepackage{amssymb}
\usepackage{booktabs}
\usepackage{graphicx}
\usepackage{caption}
\usepackage{float}
\usepackage{placeins}
\usepackage{enumitem}

\setlist{itemsep=2pt,topsep=3pt}
\newcommand{\WE}{\mathrm{WE}}
\newcommand{\FOM}{\mathrm{FOM}}
\RenewDocumentCommand{\printorcid}{}{}

\begin{document}
\restylefloat{figure}
\restylefloat{table}
\let\WriteBookmarks\relax
\def\floatpagepagefraction{1}
\def\textpagefraction{.001}

\shorttitle{Three diagnostics for neural-coordinate WE sampling claims}
\shortauthors{M.Y. Alzaatreh and H. Alsenawe}

\title[mode=title]{Three Diagnostics for Trustworthy Efficiency and Reliability
Claims in Neural-Coordinate Weighted Ensemble Sampling}

\author{M.Y. Alzaatreh}
\cormark[1]
\ead{mzaatreh@fbsu.edu.sa}
\credit{Conceptualization, Methodology, Software, Formal analysis, Investigation,
Data curation, Writing -- original draft}

\affiliation{organization={Natural Sciences Unit, Fahad Bin Sultan University},
            city={Tabuk},
            country={Saudi Arabia}}

\author{H. Alsenawe}
\credit{Validation, Writing -- review \& editing}

\affiliation{organization={Fahad Bin Sultan University},
            city={Tabuk},
            country={Saudi Arabia}}

\cortext[1]{Corresponding author}

% ---- Abstract --------------------------------------------------------------
\begin{abstract}
Reports of machine-learning-accelerated Monte Carlo methods commonly compare a
learned sampling coordinate against a hand-picked baseline using a single
point-estimate figure of merit (FOM), often computed from simulation-sweep
cost alone and from a single trained network. We show, using a Weighted
Ensemble (WE) rare-event sampler for Ising and Edwards--Anderson spin systems
as a fully worked case study, that this practice can produce headline claims
that do not survive three inexpensive diagnostic checks. First, a one-way
variance decomposition of a 17-network ensemble shows that the entire
apparent seed-to-seed spread in a published efficiency ratio is attributable
to finite-replica sampling noise, not to genuine network-initialization
variance; the between-seed variance component is exactly zero after
correction. Second, an unseeded random network initializer silently reversed
a training-objective comparison's headline result from a $+14\%$ win to a
$-17\%$ loss between two runs of an identical configuration; once seeded and
swept over 17 initializations, the comparison's own natural extensions
(finer binning, larger sampling budget) both independently fail to help the
learned coordinate, and the properly cost-corrected verdict is a significant
loss (FOM ratio $0.405$, 95\% CI $[0.268,0.623]$) rather than the
inconclusive result obtained from simulation-sweep cost alone. Third, we
summarize a companion result in which charging a learned coordinate's real
wall-clock inference cost, rather than its simulation-sweep cost, converts
an apparently favorable comparison into a $\sim$3--5$\times$ efficiency loss.
None of the three findings required new expensive computation to obtain --
each reused already-collected simulation output and cost between minutes and
a few hours of additional analysis. We package the three checks as a short,
concrete audit checklist for reporting learned-coordinate rare-event sampling
results, and argue that it should precede, not follow, any point-estimate
efficiency or reliability claim.
\end{abstract}

% ---- Highlights ------------------------------------------------------------
\begin{highlights}
\item A 17-seed variance decomposition shows an ensemble's seed-to-seed
spread is entirely finite-replica noise, not network-initialization
variance.
\item An unseeded network initializer silently reversed a headline
efficiency claim between two identical reruns.
\item Two independently tested schedule extensions (finer binning, larger
sampling budget) both fail to rescue the learned coordinate once the
seeding bug is fixed.
\item Charging real wall-clock inference cost, not simulation-sweep cost,
reverses an efficiency comparison for a companion coordinate on the same
system.
\item We propose a three-item audit checklist requiring no new expensive
simulation, only reanalysis of data already collected.
\end{highlights}

% ---- Keywords --------------------------------------------------------------
\begin{keywords}
Rare-event sampling \sep Weighted Ensemble \sep variance decomposition \sep
reproducibility \sep neural importance sampling \sep Ising model \sep spin
glass \sep figure of merit
\end{keywords}

\maketitle

% ============================================================================
\section{Introduction}\label{sec:intro}

Machine-learning-accelerated Monte Carlo methods are increasingly proposed as
replacements or supplements for hand-designed importance functions in
rare-event sampling, from neural bias potentials~\citep{hua2024} to
label-free self-consistency training~\citep{kimcai2026} to
committor-network approaches for transition-path
problems~\citep{khoo2019,li2022,kang2024,trizio2025}. A typical claim in this
literature has the shape: ``coordinate $A$ (learned) outperforms coordinate
$B$ (hand-picked) by $X\%$ on figure of merit $Y$,'' supported by one or a
handful of runs of each method.

This paper does not propose a new sampling method. Instead, using a Weighted
Ensemble (WE) rare-event sampler for Ising and Edwards--Anderson spin systems
as a fully instrumented case study, we show that three inexpensive
diagnostic checks -- none requiring new expensive simulation -- can change
the verdict of exactly this kind of claim, in one case reversing its sign
entirely. The case study is a companion project to a physics-focused paper
reporting a validated transport-inspired learned coordinate for the same
system~\citep{alzaatreh2026}; the present paper is methodological, focused on
comparisons made and abandoned during that project's development, together
with an audit of a currently used component of it.

The three diagnostics are:

\begin{enumerate}[label=(\arabic*)]
\item \textbf{Decompose seed-to-seed variance before ensembling to reduce
it.} An observed spread across independently initialized networks can be
mostly or entirely replica-sampling noise rather than genuine
initialization variance, in which case ensembling several networks pays real
inference cost to cancel a source of variance that is not there
(\S\ref{sec:diag1}).
\item \textbf{Audit random-seed control, and check that natural schedule
extensions do not silently rescue a negative result.} An unseeded
initializer can make a single comparison's sign arbitrary; once fixed, a
result should be checked against the schedule knobs a reviewer would
naturally ask about, with the mechanism of any failure identified, not just
its outcome (\S\ref{sec:diag2}).
\item \textbf{Charge a learned coordinate's real wall-clock cost, not its
simulation-sweep cost.} A cost accounting restricted to Monte Carlo sweeps is
blind to repeated network inference, which can dominate wall-clock time and
reverse an efficiency verdict (\S\ref{sec:diag3}, summarizing a companion
result).
\end{enumerate}

Each diagnostic is demonstrated on real data from the same project, using
methods (ANOVA-style variance decomposition, paired reruns, bootstrap
confidence intervals) that are standard but, in our reading of the recent
neural-importance-sampling literature, not yet routinely applied together
before a headline claim is reported. \S\ref{sec:checklist} distills the three
diagnostics into a short checklist and discusses when each is worth the
(small) additional analysis cost it requires.

% ============================================================================
\section{Case-study system}\label{sec:system}

The case study is a Weighted Ensemble (WE) sampler~\citep{huber1996,
zhang2010,zuckerman2017} targeting deep magnetization or energy tails of the
two-dimensional Ising ferromagnet and the Edwards--Anderson (EA) spin
glass~\citep{ising1925,edwards1975} on $L\times L$ lattices, with a learned
neural coordinate used to bin walkers during resampling. Two distinct
training objectives for the learned coordinate were developed over the
course of the project: a \emph{surrogate-label} objective, in which a
network is trained to regress a milestoning-style progress label collected
from pilot trajectories, and a \emph{self-consistency} objective following
Kim and Cai~\citep{kimcai2026}, in which the network is trained without
labels via a harmonic backward-equation residual. The companion physics
paper~\citep{alzaatreh2026} reports validated results for the
self-consistency objective, including a disorder-stratified reliability gain
and its associated wall-clock cost.

The present paper does not use the self-consistency results. It instead
examines two lines of analysis internal to the same project that never
appeared in the companion paper: an ensemble-variance investigation
conducted on the surrogate-label objective's 17-seed sweep
(\S\ref{sec:diag1}--\S\ref{sec:diag2}), and, briefly, the companion paper's
wall-clock cost correction viewed as a third instance of the same general
problem (\S\ref{sec:diag3}). We refer readers to~\citep{alzaatreh2026} for
the full derivation of the WE estimator, the adjoint/committor
correspondence motivating the learned coordinate, and the self-consistency
training procedure; only the details needed to follow the diagnostics below
are repeated here.

For a WE run with replicas $R$, the figure of merit used throughout is
\begin{equation}
\FOM = \frac{1}{\mathrm{rel.sd}^2 \cdot \mathrm{cost}},
\end{equation}
where $\mathrm{rel.sd}$ is the relative standard deviation of the estimated
target probability $\hat\pi$ across replicas and $\mathrm{cost}$ is a
per-replica computational cost. The central point of this paper is that
\emph{which} cost enters this denominator -- simulation sweeps only, or
total wall-clock time including network inference -- is itself a
methodological choice that can change the sign of a comparison
(\S\ref{sec:diag3}).

% ============================================================================
\section{Diagnostic 1: decompose variance before ensembling}\label{sec:diag1}

\subsection{Motivation}

A common response to noisy learned-coordinate results is to train several
independently initialized networks and pool or ensemble their outputs, on
the assumption that network-to-network variability is a meaningful source of
noise worth averaging over. This assumption is testable, and testing it is
cheap relative to the cost of training and running the extra networks.

\subsection{Method}

For $S$ independently trained networks (``seeds''), each evaluated with $R$
replicas of the downstream WE comparison, let $x_s = \log(\FOM_{I,s} /
\FOM_{E,s})$ denote the log figure-of-merit ratio of the learned coordinate
$I_\theta$ against the energy-binned baseline $E$ for seed $s$. The observed
across-seed variance,
\begin{equation}
V_{\text{across}} = \frac{1}{S-1}\sum_{s=1}^{S} (x_s - \bar x)^2,
\end{equation}
confounds two sources: genuine between-seed (initialization) variance
$V_{\text{between}}$, and within-seed variance $V_{\text{within}}$ arising
from the finite replica count $R$ used to estimate each $x_s$. A one-way
random-effects decomposition estimates the within-seed component directly,
by bootstrap-resampling each seed's own $R$ replicas independently ($10^4$
resamples per seed in this analysis), and obtains the between-seed component
by subtraction, floored at zero:
\begin{equation}
V_{\text{between}} = \max\!\big(0,\; V_{\text{across}} - \overline{V_{\text{within}}}\big).
\end{equation}

\subsection{Result}

Applied to the project's 17-seed sweep of the surrogate-label coordinate at
a fixed rarity depth ($R=10$ replicas per seed): $V_{\text{across}} =
0.2447$ ($\mathrm{sd}=0.495$), while the mean within-seed variance is
$\overline{V_{\text{within}}} = 0.3603$ ($\mathrm{sd}=0.600$) --
\emph{larger} than the raw across-seed variance itself. The subtraction
therefore floors exactly at $V_{\text{between}} = 0$ (Table~\ref{tab:anova}).
In other words, every seed's own individual measurement is already noisier
than the spread between seeds' point estimates; the apparent
seed-dependence of the published ratio is consistent with being 100\%
sampling noise, with no detectable contribution from network
initialization.

\begin{table}[h]
\centering
\caption{One-way variance decomposition of the log figure-of-merit ratio
$\log(\FOM_{I_\theta}/\FOM_E)$ across 17 independently trained networks,
$R=10$ replicas each.}
\label{tab:anova}
\begin{tabular}{lc}
\toprule
Quantity & Value \\
\midrule
Number of seeds $S$ & 17 \\
Replicas per seed $R$ & 10 \\
$V_{\text{across}}$ (raw, confounded) & 0.2447 \\
$\overline{V_{\text{within}}}$ (replica noise, bootstrapped) & 0.3603 \\
$V_{\text{between}}$ (ANOVA subtraction, floored) & 0.0000 \\
Fraction of spread attributable to replica noise & $\geq 100\%$ \\
\bottomrule
\end{tabular}
\end{table}

This has a direct practical consequence for compute allocation. A $K$-network
deep ensemble built to average over initialization variance pays $K\times$
the (non-trivial) per-network inference cost -- direct profiling of the
network used here shows a single forward pass over a WE population of
$\sim$2000 walkers costs on the order of 200~ms per WE iteration on a
single CPU core, dominated by two small convolutional layers -- to cancel a
variance component this analysis shows is not detectably present. The
statistically efficient use of additional compute here is more replicas per
existing seed, not more seeds (Table~\ref{tab:allocation}), though we note
replica cost for a learned coordinate is not free either, since each
additional replica still pays the network's per-iteration inference cost.

\begin{table}[h]
\centering
\caption{Predicted standard error of the pooled mean log-ratio under two
compute-allocation strategies, from the fitted ANOVA model of
Table~\ref{tab:anova}. The current configuration ($S=17$, $R=10$) does not
resolve the ratio from unity; both strategies can, at different total
compute.}
\label{tab:allocation}
\begin{tabular}{lcc}
\toprule
Strategy & Allocation & Predicted SE \\
\midrule
More seeds, fixed $R=10$ & $S=17$ (current) & 0.146 \\
 & $S=50$ & 0.085 \\
 & $S=150$ & 0.049 \\
More replicas, fixed $S=17$ & $R=10$ (current) & 0.146 \\
 & $R=80$ & 0.052 \\
 & $R=320$ & 0.026 \\
\bottomrule
\end{tabular}
\end{table}

\FloatBarrier

% ============================================================================
\section{Diagnostic 2: seed control and schedule-extension audits}\label{sec:diag2}

\subsection{A silent sign reversal}

The 17-seed sweep behind \S\ref{sec:diag1} exists because the project's
first reported result for this comparison -- a single run showing the
learned coordinate winning by $+14\%$ ($\FOM_{I_\theta}=4.811\times10^{-8}$
vs.\ $\FOM_E = 4.215\times10^{-8}$) -- did not reproduce. A second run of the
byte-identical configuration (same cached training labels, same
hyperparameters) gave $\FOM_{I_\theta} = 3.484\times10^{-8}$, a $-17\%$
\emph{loss} against the same $\FOM_E$ baseline (Table~\ref{tab:seedbug}).
The two runs' naive-sampling and energy-baseline figures matched exactly
across both runs; only the trained-network-dependent figure differed. The
root cause was a missing \texttt{torch.manual\_seed()} call: the network's
weight initialization was left to the platform's default (unseeded) random
state, so every rerun trained a different network from the same nominal
configuration.

\begin{table}[h]
\centering
\caption{Two runs of an identical configuration, before the seeding fix.
naive and $\WE[E]$ (energy-binned, no learned network) are bit-identical
across runs; only $\WE[I_\theta]$ (learned network) differs.}
\label{tab:seedbug}
\begin{tabular}{lcc}
\toprule
Method & Run 1 FOM & Run 2 FOM \\
\midrule
naive & $1.476\times10^{-7}$ & $1.476\times10^{-7}$ \\
$\WE[m]$ & $1.173\times10^{-7}$ & $1.173\times10^{-7}$ \\
$\WE[E]$ & $4.215\times10^{-8}$ & $4.215\times10^{-8}$ \\
$\WE[I_\theta]$ & $4.811\times10^{-8}$ ($+14\%$) & $3.484\times10^{-8}$ ($-17\%$) \\
\bottomrule
\end{tabular}
\end{table}

This is a narrow, mechanical bug, but its consequence is not narrow: absent
the second run, the project would have reported a clean single-comparison
win as its headline learned-coordinate result. After fixing the seed and
sweeping 17 independent initializations, point-ratio estimates for
individual seeds range from $-52\%$ to $+112\%$ around unity -- consistent
with the $+14\%$ and $-17\%$ single-run outcomes both being unremarkable
draws from the same wide distribution, not evidence of a real effect in
either direction from any one run.

\subsection{Two schedule extensions independently fail to help}

A natural response to a wide, centered-near-unity comparison is to ask
whether a finer discretization or a larger sampling budget would let the
learned coordinate's advantage (if any) emerge more clearly. Two such
extensions were tested at a deeper rarity target (``beyond=2''), each
against the same baselines, and each independently failed:

\begin{itemize}
\item \textbf{Finer binning} ($n_{\text{bins}}$ from 50 to 90, replicas
fixed at 10): the learned coordinate's own relative standard deviation
improved only marginally (0.416 $\to$ 0.395) while its cost nearly doubled
($6.90\times10^{8} \to 1.19\times10^{9}$), so its own FOM \emph{worsened}
($8.360\times10^{-9} \to 5.363\times10^{-9}$). The energy baseline's cost is
structurally insensitive to bin count (it stayed exactly
$5.45\times10^{8}$ across both configurations), because energy is a
discrete, quantized coordinate on this lattice that saturates a small fixed
number of occupied bins regardless of $n_{\text{bins}}$ -- extra bins can
only add cost to the continuous learned coordinate, not help it.
\item \textbf{Larger sampling budget} (WE iterations, $\mathtt{we\_iter}$,
1000 to 2000, $n_{\text{bins}}=50$): cost scaled worse than the expected
$2\times$ for every method (population grows as the run pushes further into
the tail), and every method's FOM \emph{worsened}, the learned coordinate's
most severely
($\FOM_{I_\theta}: 8.360\times10^{-9} \to 1.698\times10^{-9}$, an 80\%
drop; its own relative standard deviation \emph{increased}, 0.416 $\to$
0.540, despite $3\times$ the sampling).
\end{itemize}

Table~\ref{tab:deadends} summarizes both. Two structurally different
extensions, tested independently, produce the same qualitative outcome:
neither rescues the learned coordinate at this depth, and identifying the
cost mechanism in each case (a saturating discrete coordinate for binning; a
population that grows faster than expected with iteration budget) rules out
re-attempting either lever as a productive next step, rather than leaving
the question open for a future, more expensive retry.

\begin{table}[h]
\centering
\caption{Two schedule extensions at a fixed deeper rarity target
(``beyond=2''), both against the same fixed baselines. Both independently
worsen the learned coordinate's own FOM relative to its own prior value.}
\label{tab:deadends}
\begin{tabular}{lccc}
\toprule
Extension & Baseline & Extended & Change in $\FOM_{I_\theta}$ \\
\midrule
$n_{\text{bins}}$: 50 $\to$ 90 & $8.360\times10^{-9}$ & $5.363\times10^{-9}$ & $-36\%$ \\
$\mathtt{we\_iter}$: 1000 $\to$ 2000 & $8.360\times10^{-9}$ & $1.698\times10^{-9}$ & $-80\%$ \\
\bottomrule
\end{tabular}
\end{table}

\subsection{Cost-corrected verdict}

Reporting a training-objective comparison's cost from simulation sweeps
alone, as in Table~\ref{tab:seedbug} and the FOM values above, is itself
subject to the same blindness raised in \S\ref{sec:diag3}: the learned
coordinate's repeated network inference is uncounted. Recomputing the
17-seed sweep's pooled figure of merit using measured wall-clock time
(mean $\WE[E]$ native-cost wall time: 397~s over 16 reconfirmed seeds)
rather than simulation-sweep cost, with a two-level bootstrap resampling
both the seed axis and each seed's replica axis, gives a median ratio of
\textbf{0.405} (95\% CI $[0.268, 0.623]$, excluding unity) -- a
statistically significant loss for the learned coordinate, in contrast to
the simulation-sweep-cost figure, whose 95\% CI ($[0.84, 1.75]$ in the
project's original analysis) includes unity and would be reported as
inconclusive.

% ============================================================================
\section{Diagnostic 3: charge real inference cost, not sweep cost}\label{sec:diag3}

The two diagnostics above already used wall-clock-corrected costs in their
final reported numbers. This section states the general point explicitly
and summarizes independent confirmation from the companion
paper~\citep{alzaatreh2026}, which uses a different training objective
(self-consistency, not surrogate-label) on the same system.

A cost field computed as (number of Monte Carlo sweeps) $\times$ (per-sweep
cost) is, by construction, blind to any computation that scales with WE
\emph{iterations} rather than sweeps -- chiefly, repeated forward passes of
a trained network used for binning. Direct instrumented profiling confirms
this cost is real and not a batching or threading artifact: the network's
forward pass is already correctly batched (one call per WE iteration over
the full population, gradient tracking disabled, single-threaded as
intended), and its cost is genuine compute-bound CPU time, dominated by two
small convolutional layers, that scales with population size.

Measured directly from production run wall-clock timestamps, the learned
coordinate's WE stage consistently costs 3--5$\times$ the energy baseline's
wall-clock time across every configuration tested in this project family
(range 2.87--4.97$\times$ across more than twenty independent seed and
disorder combinations, both training objectives). In the companion
paper's self-consistency-based headline result, charging this real cost
converts the comparison from favorable-looking on a count-based reliability
metric (which is cost-independent and unaffected) to unfavorable on a
cost-normalized efficiency metric: a wall-clock-fair figure-of-merit ratio
of $0.33$ (95\% CI $[0.19,0.58]$)~\citep{alzaatreh2026}. The direction and
approximate magnitude of this correction closely matches the surrogate-label
objective's independently derived $0.405\times$ figure in
\S\ref{sec:diag2}, despite the two training objectives, and the two
diagnostics that produced them, being otherwise unrelated -- consistent
with cost-blindness being a property of the cost accounting, not of either
specific learned coordinate.

% ============================================================================
\section{An audit checklist}\label{sec:checklist}

The three diagnostics above share a structure: each is a specific,
inexpensive reanalysis of data already collected, not a request for more
expensive simulation, and each is capable of overturning a headline
efficiency or reliability claim on its own. We summarize them as a short
checklist for reporting learned-coordinate rare-event sampling results.

\begin{enumerate}[label=(\arabic*)]
\item \textbf{Before ensembling networks to reduce variance, decompose the
variance you are trying to reduce.} If a one-way (or richer) variance
decomposition shows the between-seed component is small or zero relative to
per-seed sampling noise, the ensemble's inference cost buys nothing, and
additional replicas of existing networks are the more efficient use of the
same compute (\S\ref{sec:diag1}).
\item \textbf{Seed every stochastic component of a trained-coordinate
pipeline, and confirm determinism by rerunning an identical configuration
at least once before reporting a single-run result as a finding.} A cheap
rerun is the entire cost of catching a bug capable of reversing a headline
sign (\S\ref{sec:diag2}). Once seeded, test the schedule extensions a
reviewer would naturally suggest (finer discretization, larger sampling
budget); if they fail, identify the cost mechanism responsible rather than
leaving the failure unexplained, since the mechanism determines whether the
lever is worth retrying at a different scale.
\item \textbf{Compute the reported figure of merit's cost term from
measured wall-clock time for every method being compared, not from a
simulation-sweep count that is blind to learned-component inference.} Where
wall-clock measurement is not yet available, treat a sweep-cost-only
comparison as provisional and flag it as such, since the correction is not
guaranteed to be small (\S\ref{sec:diag3}).
\end{enumerate}

None of the three checks requires new expensive simulation on the scale
of the original comparison being audited: the variance decomposition
(\S\ref{sec:diag1}) reuses already-collected replica data; the seed audit
(\S\ref{sec:diag2}) requires one confirmatory rerun of an existing
configuration, cheap relative to a full sweep; and the cost correction
(\S\ref{sec:diag3}) requires only that wall-clock time already be recorded
(or, if not, added as a field going forward) rather than any new
simulation. We suggest this ordering -- variance decomposition, seed and
schedule audit, cost correction -- as a reasonable default sequence, since
each check is comparably cheap and any one of them failing can make the
others' outcome moot for the specific claim at hand.

We do not claim these three are exhaustive, nor that every
learned-coordinate rare-event sampling claim in the literature is subject to
all three problems; we report only what a systematic reanalysis of one
project's own internal history found. The consistency of direction across
two independently developed training objectives on the same underlying
system (\S\ref{sec:diag3}) is, however, some evidence that the cost-blindness
problem at least is not an artifact specific to one coordinate design.

% ============================================================================
\section{Conclusions}\label{sec:conclusions}

Using a fully instrumented Weighted Ensemble rare-event sampler for Ising
and Edwards--Anderson spin systems, we showed that three inexpensive
diagnostic checks -- a variance decomposition, a seed and schedule audit,
and a wall-clock cost correction -- each independently change the reported
verdict of a learned-coordinate efficiency or reliability comparison drawn
from the same project. A 17-network ensemble's entire seed-to-seed spread
was shown to be finite-replica noise rather than genuine initialization
variance. An unseeded initializer silently reversed a headline result
between two runs of an identical configuration, and once fixed, two natural
schedule extensions were shown, with identified mechanism, to independently
fail to rescue the comparison. Charging real wall-clock inference cost
rather than simulation-sweep cost reversed an efficiency verdict twice, for
two different training objectives on the same system, from an inconclusive
or favorable-looking result to a statistically significant loss. None of
these findings required new expensive computation; each reused data or
compute already paid for by the underlying project. We offer the three
checks as a short, concrete checklist that we believe should be applied
before, not after, a learned-coordinate rare-event sampling paper's headline
efficiency or reliability claim is finalized.

% ============================================================================
\section*{Data and Code Availability}
The scripts implementing the three diagnostics (variance decomposition,
seed-sweep and schedule-extension drivers, wall-clock bootstrap correction),
along with the underlying run logs and result files referenced in
Tables~\ref{tab:anova}--\ref{tab:deadends}, are available at
\url{https://github.com/mzaatreh-phs/ising-transport-rare-event}, the same
repository as the companion paper~\citep{alzaatreh2026}.
% Optional author action before submission: add a versioned archival DOI here
% after depositing the repository in Zenodo.

\section*{Acknowledgements}
The authors gratefully acknowledge the support of Fahad Bin Sultan
University, Tabuk, Saudi Arabia, including access to the computational
environment used for the numerical experiments.

\printcredits

\section*{Declaration of competing interest}
The authors declare no known competing financial interests or personal
relationships that could have appeared to influence the work reported in
this paper.

\section*{Declaration of generative AI and AI-assisted technologies in the manuscript preparation process}
During the preparation of this work, the authors used Claude (Anthropic) to
assist with manuscript organization, language refinement, and software
support. After using this tool, the authors reviewed and edited the content
as needed, independently verified the scientific claims and numerical
results, and take full responsibility for the content of the publication.

% ============================================================================
\begin{thebibliography}{99}
\bibitem{huber1996} G. A. Huber, S. Kim, Weighted-ensemble Brownian dynamics
simulations for protein association reactions, Biophys. J. 70 (1996)
97--110.
\bibitem{zhang2010} B. W. Zhang, D. Jasnow, D. M. Zuckerman, The ``weighted
ensemble'' path sampling method is statistically exact for a broad class of
stochastic processes and binning procedures, J. Chem. Phys. 132 (2010)
054107.
\bibitem{zuckerman2017} D. M. Zuckerman, L. T. Chong, Weighted ensemble
simulation: review of methodology, applications, and software, Annu. Rev.
Biophys. 46 (2017) 43--57.
\bibitem{ising1925} E. Ising, Beitrag zur Theorie des Ferromagnetismus, Z.
Phys. 31 (1925) 253--258.
\bibitem{edwards1975} S. F. Edwards, P. W. Anderson, Theory of spin glasses,
J. Phys. F: Met. Phys. 5 (1975) 965--974.
\bibitem{hua2024} X. Hua, R. Ahmad, J. Blanchet, W. Cai, Accelerated
sampling of rare events using a neural network bias potential,
arXiv:2401.06936 (2024).
\bibitem{kimcai2026} M. Kim, W. Cai, Accelerated Markov Chain Monte Carlo
simulation via neural network-driven importance sampling, arXiv:2602.12294
(2026).
\bibitem{khoo2019} Y. Khoo, J. Lu, L. Ying, Solving for high-dimensional
committor functions using artificial neural networks, Res. Math. Sci. 6
(2019) 1.
\bibitem{li2022} H. Li, Y. Khoo, Y. Ren, L. Ying, A semigroup method for
high dimensional committor functions based on neural network, Proc. Mach.
Learn. Res. 145 (2022) 598--618.
\bibitem{kang2024} P. Kang, E. Trizio, M. Parrinello, Computing the
committor with the committor to study the transition state ensemble, Nat.
Comput. Sci. 4 (2024) 451--460.
\bibitem{trizio2025} E. Trizio, P. Kang, M. Parrinello, Everything
everywhere all at once: a probability-based enhanced sampling approach to
rare events, Nat. Comput. Sci. 5 (2025) 582--591.
\bibitem{efron1979} B. Efron, Bootstrap methods: another look at the
jackknife, Ann. Statist. 7 (1979) 1--26.
\bibitem{alzaatreh2026} M. Y. Alzaatreh, Transport-Inspired Rare-Event
Sampling in Spin Systems: Learned Importance Improves Deep-Target
Reliability, submitted (2026). % TODO(authors): update with journal name,
% volume, and DOI once available; this is the companion paper referenced
% throughout \S\ref{sec:system}--\S\ref{sec:diag3}.
\end{thebibliography}

\end{document}
